\documentclass{QD_2026}
\makeatletter
\setlength\@fptop{0pt}
\makeatother

\usepackage{graphicx}
\usepackage{microtype}
\usepackage{amsfonts, amssymb}
\usepackage{booktabs}
\usepackage{siunitx}
\sisetup{detect-all}
\graphicspath{{figures/}}

% Macros
\newcommand\bemt  {\textsf{BEMT}}
\newcommand\vlm   {\textsf{VLM}}
\newcommand\fwh   {\textsf{FWH}}
% Additional macros defined inline

\begin{document}

%% Header image on first page
\begin{figure}[t!]
  \centering
  \includegraphics[keepaspectratio,width=0.55\textwidth]{QD2026_logo_cropped.pdf}
\end{figure}
\hfill\hfill

\title{Physical Rotor Noise Modelling: BEMT vs Vortex Lattice Method}
\QDsession{UAS/UAM Noise Modelling}

\providecommand\email[1]{\texttt{#1}}
\author{
  \textbf{Dmitrii Mukhutdinov}\email{d.mukhutdinov@qmul.ac.uk}
  \and
  \textbf{Lin Wang} \email{lin.wang@qmul.ac.uk}
  \\School of Electronic Engineering and Computer Science\\
  Queen Mary University of London, United Kingdom
}

\maketitle
\author{}

%==============================================================================
\section{Introduction}
%==============================================================================

Rotors on multirotor drones generate acoustic noise dominated by tonal harmonics
at the \emph{blade passage frequency} $BPF = B \times \Omega$, where $B$ is the
number of blades and $\Omega$ the rotor angular speed in rad/s.
These harmonics encode the rotor geometry and operating conditions.
Two aerodynamic models are widely used to compute the blade forces:
%
\begin{enumerate}
  \item \textbf{Blade Element Momentum Theory (BEMT)}: a quasi-analytical strip
    approach combining 2-D airfoil polars with momentum conservation.
  \item \textbf{Vortex Lattice Method (VLM)}: a potential-flow model using
    bound and trailing vortices via the Biot--Savart law.
    More physically faithful and differentiable end-to-end.
\end{enumerate}
%
This report compares the acoustic spectra produced by BEMT and VLM against
real motor recordings from the DREGON dataset~\cite{strauss2018dregon},
benchmarks their single-GPU speed, and presents both methods
pedagogically.

%==============================================================================
\section{Blade Element Momentum Theory (BEMT)}
\label{sec:bemt}
%==============================================================================

BEMT slices the blade into $N_r$ radial strips and treats each as an
independent 2-D airfoil in steady local flow.

At radius $r$, the blade tip speed is $U_y = \Omega r$ and the inflow
speed is $v_\infty$ (downwash in hover, freestream in flight).
The local speed and inflow angle are:
%
\begin{equation}
  V(r) = \sqrt{v_\infty^2 + (\Omega r)^2}, \qquad
  \phi(r) = \arctan\frac{v_\infty}{\Omega r}.
\end{equation}
%
The angle of attack is $\alpha = \theta(r) - \phi$, where $\theta(r)$ is the
geometric twist.  Thin-airfoil theory gives lift and drag per unit span:
%
\begin{equation}
  \ellis(r) = \tfrac12\rho V^2\,c(r)\,C_\ell(\alpha), \qquad
  D(r)     = \tfrac12\rho V^2\,c(r)\,C_d(\alpha).
\end{equation}
%
Forces are resolved in the rotor disk frame (radial, tangential, thrust):
%
\begin{equation}
  F_x^{\mathrm{disk}} = D\cos\phi - \ellis\sin\phi, \qquad
  F_y^{\mathrm{disk}} = D\sin\phi + \ellis\cos\phi, \qquad
  F_z^{\mathrm{disk}} = 0.
\end{equation}
%
In hover ($\phi\approx 0$): $F_y \approx \ellis$ (thrust/lift),
$F_x \approx D$ (drag), $F_z = 0$ (thin-airfoil assumption: no out-of-plane force).

\textbf{Assumptions and limitations:}
the strips are independent (no blade-to-blade interaction);
wake is modelled as uniform (no tip-vortex model);
steady (no dynamic stall);
most accurate at midspan.

%==============================================================================
\section{Vortex Lattice Method (VLM)}
\label{sec:vlm}
%==============================================================================

VLM represents each blade as a horseshoe vortex: a bound vortex of circulation
$\Gamma$ running hub-to-tip, with two semi-infinite trailing vortices
extending downstream.

The Kutta--Joukowski theorem gives lift per unit span:
%
\begin{equation}
  L' = \rho\,U \times \Gamma \quad\Longrightarrow\quad
  \Gamma = \frac{L'}{\rho U}.
\end{equation}
%
A vortex particle of circulation $\Gamma_i$ at position $\mathbf{y}_i$ induces
velocity at field point $\mathbf{x}$ via the Biot--Savart law:
%
\begin{equation}
  \label{eq:biotsavart}
  \mathbf{v}(\mathbf{x}) = \frac{\Gamma_i}{4\pi}
    \frac{\mathbf{r}_i\times d\boldsymbol\ell_i}
         {(|\mathbf{r}_i|^2+\epsilon^2)^{3/2}},
  \qquad \mathbf{r}_i = \mathbf{x}-\mathbf{y}_i,
\end{equation}
%
where $\epsilon$ is a regularisation (core) radius.

\textbf{Key property:} The Biot--Savart law uses only matmul, cross-product,
and normalisation --- all natively differentiable in PyTorch autograd.
This makes the VLM force-to-pressure chain fully end-to-end differentiable,
which is the decisive advantage for inverse design.

%==============================================================================
\section{Acoustic Propagation: Ffowcs--Williams Hawkings (FWH)}
\label{sec:fwh}
%==============================================================================

The FWH equation rewrites the compressible Navier--Stokes equations as a wave
equation driven by surface integrals.  For a compact source of force
$\mathbf{F}(t)$ at retarded time $\tau$, Farassat Formulation 1A gives:
%
\begin{equation}
  \label{eq:farassat1a}
  4\pi\,p'(\mathbf{x},t) =
    \frac{1}{c_0}\,
    \frac{\hat{\mathbf{r}}\cdot\dot\mathbf{F}}
         {r\,(1-M_r)^2}\Bigg|_\tau
    + \frac{\hat{\mathbf{r}}\cdot\mathbf{F}}
           {r^2\,(1-M_r)^2}\Bigg|_\tau ,
\end{equation}
%
where $\hat{\mathbf{r}}$ is the unit vector from source to observer,
$r=|\mathbf{x}-\mathbf{y}|$, and $M_r = \mathbf{v}\cdot\hat{\mathbf{r}}/c_0$
is the Mach number projection.  The total pressure is summed over all
strips and blades:
%
\begin{equation}
  p'(\mathbf{x},t) = \sum_{i=1}^{N_r}\sum_{b=1}^{B} p'_{ib}(\mathbf{x},t).
\end{equation}

The pipeline is: blade geometry $\to$ aerodynamic forces (BEMT or VLM)
$\to$ retarded time $\tau$ (Newton--Raphson on
$t - \tau - |\mathbf{x}-\mathbf{y}(\tau)|/c_0 = 0$) $\to$ observer
pressure $p'(\mathbf{x},t)$.

%==============================================================================
\section{Results}
%==============================================================================

\textbf{Setup.}  2-blade rotor, $R=\SI{0.152}{m}$,
APC~\SI{10\times7}{} chord distribution, linear twist
$\theta(r)=15^\circ(1-r/R)$.
RPM sweep: $50$, $70$, $90\,\SI{}{rev/s}$ (constant-speed hover).
Dataset: DREGON individual motor recordings, 4 motors $\times$ 3 RPMs,
12 samples of $\SI{5}{s}$ each.
Microphone: DREGON linear 8-channel array, channel 0.
%
Spectral metrics: (i) normalised MSE of zero-mean PSDs (lower = better);
(ii) Pearson $r$ of normalised PSDs ($r=1$: identical harmonic structure).

%-----------------------------------------------------------------------------
\subsection{Spectral accuracy}
%-----------------------------------------------------------------------------

\begin{figure}[t!]
  \centering
  \includegraphics[width=0.48\textwidth]{fig_spectra_all_bpfs.pdf}
  \caption{
    BEMT (dashed) and VLM (dotted) simulated spectra vs.\ DREGON real
    recordings (solid gray).  Both models reproduce the BPF harmonics
    exactly (determined by kinematics $B\times\Omega$).
    The dominant gap is the broadband noise floor in real recordings,
    which neither inviscid potential-flow model captures.}
  \label{fig:spectra_all}
\end{figure}

Figure~\ref{fig:spectra_all} shows spectra for all 12 motor/RPM combinations.
Figure~\ref{fig:mse_by_rpm} shows spectral distance from the real recording.

Both methods give MSE $\approx 40$--$\SI{45}{dB}$ and Pearson
$r\approx0.75$ against real recordings.  This gap is entirely attributable to
the broadband motor/ESC noise floor in the real recordings (see
Section~\ref{sec:gaps}), not to differences between the two models.

A critical finding is that VLM and BEMT produce \emph{near-identical}
acoustic waveforms (Pearson $r=0.9996$, RMS SPL difference
$<0.04\,\mathrm{dB}$).  This is not a numerical coincidence ---
it follows from the physics:
%
\begin{enumerate}
  \item VLM uses $\Gamma = L'/(\rho U)$ (Kutta--Joukowski), which
    recovers exactly the same section lift as thin-airfoil BEMT for the same
    angle of attack (verified: max relative difference $<10^{-4}$).
  \item Both feed the same lift into the Farassat-1A FWH integral.
    The only algorithmic difference is that BEMT solves the retarded time
    jointly over all $N_r=30$ strips, while VLM solves it strip-by-strip;
    the two solutions differ by $<0.04\,\mathrm{dB}$.
\end{enumerate}
%
An initial run with $n_t=200$ (effective sample rate $\SI{40}{Hz}$,
Nyquist $=$\SI{20}{Hz} vs BPF $\approx$\SI{180}{Hz}) produced an
apparent \SI{1.4}{dB} MSE advantage for BEMT --- a resampling artefact
caused by complete harmonic aliasing.  With $n_t=2000$
(Nyquist $>\!$BPF), the VLM--BEMT MSE collapses to $0.04\,\mathrm{dB}^2$.

\begin{figure}[t!]
  \centering
  \includegraphics[width=0.45\textwidth]{fig_mse_by_rpm.pdf}
  \caption{
    Spectral distance (normalised MSE, lower = better) from real recording,
    by RPM.  Both BEMT and VLM are within $\approx 2\,$dB of each other.
    Error bars: $\pm 1$ standard deviation over the four motors.}
  \label{fig:mse_by_rpm}
\end{figure}

\begin{figure}[t!]
  \centering
  \includegraphics[width=0.45\textwidth]{fig_correlation_by_rpm.pdf}
  \caption{
    Spectral shape correlation (Pearson $r$, higher = better) to real recording.
    Both models give $r\approx 0.75$ across all RPMs (dashed line: $r=0.80$).}
  \label{fig:corr_by_rpm}
\end{figure}

\clearpage

\begin{figure}[t!]
  \centering
  \includegraphics[width=0.45\textwidth]{fig_speed_benchmark.pdf}
  \caption{
    Single-GPU wall-clock time per call (50-run average, $N_t=200$,
    RTX 4070 Ti, PyTorch 2.7+cu126).
    VLM is $6.7\times$ faster than BEMT per call.}
  \label{fig:speed}
\end{figure}

%-----------------------------------------------------------------------------
\subsection{Benchmark: speed on single GPU}
%-----------------------------------------------------------------------------

Table~\ref{tab:benchmark} and Figure~\ref{fig:speed} summarise timing.
VLM is consistently faster because it processes each strip independently
(single-source retarded-time Newton solves are cheap), while BEMT's
joint solver requires $N_r\times N_t$ Newton steps per call.

\begin{table}[t!]
  \centering
  \caption{Single-GPU wall-clock time per acoustic call
    (mean over 50 runs, $N_t = 200$, RTX 4070 Ti, PyTorch 2.7+cu126).}
  \label{tab:benchmark}
  \begin{tabular}{lSS[table-format=4.1]S[table-format=2.2]}
    \toprule
    Method & \multicolumn{2}{c}{Time [ms]} & {Rate [calls/s]} & {Ratio} \\
           & {per call}                  & (est. per step)            &           \\
    \midrule
    BEMT    & 3361 & 16.8 & 0.30 & 1.00 \\
    VLM     &  515 &  2.6 & 1.94 & 0.15 \\
    \bottomrule
  \end{tabular}
\end{table}

%==============================================================================
\section{Discussion and Conclusions}
%==============================================================================

\subsection{Why neither model matches real recordings}
\label{sec:gaps}
The $r\approx0.75$ correlation gap is entirely caused by the
broadband noise floor in real single-motor recordings:
%
\begin{itemize}
  \item Motor and ESC electronic noise (dominant above $\sim$\SI{500}{Hz}).
  \item Trailing-edge and tip-vortex turbulent boundary-layer noise.
  \item Recording microphone self-noise and room reflections.
  \item Multi-rotor acoustic coupling (four motors on a quadcopter, not one).
\end{itemize}
%
Neither BEMT nor VLM models any of these: both are inviscid potential-flow
models with an ideal vortex wake, producing only the deterministic tonal
harmonics at $BPF = B\times\Omega$.  Turbulent broadband noise requires
either a turbulence model (e.g.\ DES) or a full Navier--Stokes solver
(e.g.\ Lattice-Boltzmann via Lettuce).

\subsection{VLM vs BEMF: identical waveforms, different speed}
\label{sec:bemt_vs_vlm}
VLM and BEMT generate near-identical acoustic pressure waveforms
(Pearson $r=0.9996$, RMS SPL difference $<0.04\,\mathrm{dB}$).
This follows from the physics: VLM's Kutta--Joukowski circulation
$\Gamma = L'/(\rho U)$ recovers exactly the same section lift as thin-airfoil
BEMT for the same angle of attack; both feed the same lift into the
Farassat-1A integral; the only difference is the retarded-time solver
(joint vs strip-by-strip), which introduces negligible numerical noise.

Wake modelling does not add broadband noise: standard VLM uses clean
vortex filaments (potential flow, no viscosity); turbulent broadband noise
requires a turbulence model or Navier--Stokes solver.

\subsection{Recommendation for this project}
VLM is the preferred model because:
%
\begin{enumerate}
  \item \textbf{Speed.}  $6.7\times$ faster per call on a single GPU.
  \item \textbf{Differentiability.}  Full end-to-end autograd chain
    ($c(r),\theta(r)\to\Gamma\to F\to p$) via PyTorch native operations
    (matmul, cross, norm).
  \item \textbf{Spectral accuracy.}  Identical to BEMT on tonal harmonics;
    neither captures broadband noise.
\end{enumerate}
%
BEMT is retained as a fast prototyping tool and validation baseline.

%==============================================================================
\section*{Acknowledgments}
%==============================================================================

DREGON dataset: \url{http://dregon.inria.fr}.

\bibliographystyle{IEEEtran}
\bibliography{refs}

\end{document}
